\documentclass[12pt]{article}

\usepackage[a4paper,margin=2.5cm]{geometry}

\usepackage[onehalfspacing]{setspace}

\usepackage{amsmath}
\usepackage{amssymb}
\usepackage{xcolor}
\usepackage{mathtools}
\usepackage{amsthm}
\usepackage{tikz-cd}
\newtheorem{definition}{Definition}
\newtheorem{example}{Example}

\newenvironment{solution}
{\begin{proof}[Solution]}
{\end{proof}}

\newenvironment{Remark}
{\begin{proof}[Remark]}
{\end{proof}}

\newenvironment{Goal}
{\begin{proof}[Goal]}
{\end{proof}}

\setlength{\parindent}{0pt}

\begin{document}

\begin{titlepage}
    \centering

    \vspace*{\fill}

    {\Huge Linear Transformations III}

    \vspace{1cm}

    {\Large Jin-Ren Ryan Wang\par}

    \vspace{0.5cm}

    {\large August 1, 2026\par}

    \vspace*{\fill}
\end{titlepage}

\tableofcontents

\newpage

\section{Row operations and elementary row matrices}

Historically, matrices were invented and used to simplify the process of
solving systems of linear equations. One key technique for resolving such
systems is the method of Gaussian elimination, which involves three types of
fundamental operations on a matrix $A$:

\begin{enumerate}
    \item[(I)] Interchanging two rows of $A$.
    \item[(II)] Multiplying any row of $A$ by a nonzero scalar.
    \item[(III)] Adding any scalar multiple of a row of $A$ to another row.
\end{enumerate}

These operations are referred to as elementary row operations. Our goal is to
develop a concise mathematical theory that underlies these operations. To
begin with, we define three types of elementary matrices.
\subsection{Definition}

\textbf{Definition 6.1.1.}
An $(m \times m)$ elementary row matrix is a matrix obtained by
performing exactly one elementary row operation on the
$(m \times m)$ identity matrix $\mathbf{I}_m$.

We denote the three types of elementary row matrices by
\[
E^{\mathrm{I}}_{i,j},
\quad
E^{\mathrm{II}}_{\lambda,i},
\quad
E^{\mathrm{III}}_{\lambda,i,j},
\text{ respectively.}
\]

\textbf{Definition 6.1.2.}
An $m\times n$ matrix $A$ is said to be in row echelon form (REF)
if the following conditions are satisfied.

\begin{enumerate}
    \item Every nonzero row is above every zero row.
    \item The first nonzero entry (leading entry) of each nonzero row
    occurs strictly to the right of the first nonzero entry of the
    previous row.
\end{enumerate}

\subsection{Theorem}

\textbf{Theorem 6.1.1.}
Let $A$ be a matrix with $m$ rows. Then the matrices obtained by
\begin{enumerate}
    \item performing one of the row operations (I), (II), or (III) on $A$, and
    \item multiplying the corresponding elementary row matrix $E$ to the left of $A$, i.e. the matrix
    $EA$
\end{enumerate}
are the same.

\begin{proof}
The result follows directly from the definition of matrix multiplication.
\end{proof}

\textbf{Theorem 6.1.2.}
Every (square) elementary row matrix is invertible, and its inverse is given by
\[
\left(E^{\mathrm{I}}_{i,j}\right)^{-1}=E^{\mathrm{I}}_{i,j},\qquad
\left(E^{\mathrm{II}}_{\lambda,i}\right)^{-1}=E^{\mathrm{II}}_{\frac{1}{\lambda},i},\qquad
\left(E^{\mathrm{III}}_{\lambda,i,j}\right)^{-1}=E^{\mathrm{III}}_{-\lambda,i,j},
\text{ respectively.}
\]

\begin{proof}
For each type of elementary row operation:

\begin{enumerate}
    \item[(I)] Interchange rows $i$ and $j$ of $A$.
    \item[(II)] Multiply the $i$-th row of $A$ by a nonzero scalar \fcolorbox{red}{white}{$\lambda$}.
    \item[(III)] Add \fcolorbox{red}{white}{$\lambda$} times the $i$-th row of $A$ to the $j$-th row.
\end{enumerate}

Each operation can be undone by the corresponding inverse row operation:

\begin{enumerate}
    \item[(I)] Interchange rows $i$ and $j$ again.
    \item[(II)] Multiply the $i$-th row by $\dfrac{1}{\lambda}$.
    \item[(III)] Add \fcolorbox{red}{white}{$-\lambda$} times the $i$-th row to the $j$-th row.
\end{enumerate}

These inverse row operations correspond to the elementary matrices $E^{\mathrm{I}}_{i,j}$,
$E^{\mathrm{II}}_{\frac{1}{\lambda},i}$, $E^{\mathrm{III}}_{-\lambda,i,j}$
respectively. Hence every elementary row matrix is invertible with the stated inverse.
\end{proof}

\textbf{Theorem 6.1.3.}

Let $A$ be a matrix with $m$ rows. If $Q$ is an invertible
$m \times m$ matrix \textcolor{red}{\text{(i.e. $Q^{-1}$ exists)}}, then
\[
\operatorname{rank}(QA)=\operatorname{rank}(A).
\]
In particular, left multiplication by an invertible matrix preserves the rank.

\begin{proof}
By the rank inequality,
\[
\left\{
\begin{aligned}
\operatorname{rank}(QA) &\le \operatorname{rank}(A)\\
\operatorname{rank}(A) &= \operatorname{rank}(Q^{-1}QA) \le \operatorname{rank}(QA)
\end{aligned}
\right.  
\]

Since $Q$ is invertible,
\[
A=Q^{-1}(QA),
\]
and hence
\[
\operatorname{rank}(A)=\operatorname{rank}\!\left(Q^{-1}(QA)\right)\le \operatorname{rank}(QA).
\]

Therefore,
\[
\operatorname{rank}(QA)=\operatorname{rank}(A).
\]

Since \textcolor{red}{all elementary row matrices E are invertible}, it follows that
\[
\operatorname{rank}(\textcolor{red}{E}A)=\operatorname{rank}(A)
\]
for every elementary row matrix $E$. Hence elementary row operations preserve the rank.

\begin{Goal}
To compute the rank of a matrix $A$, we apply elementary row operations
to row-reduce $A$ into a simpler form. Since elementary row operations
preserve the rank, the rank of the reduced matrix is equal to the rank
of the original matrix.
\end{Goal}

\end{proof}

\textbf{Theorem 6.1.4.}
If $A$ is in its REF, then $\operatorname{rank}(A)=\text{the number of nonzero rows of }A$.

\begin{proof}
We postpone the proof of this result. (Although it is easy to verify
the statement by considering concrete examples.)
\end{proof}

\textbf{Theorem 6.1.5.}
Let $A$ be a matrix with $m$ rows. $A$ can be transformed into a REF by a finite number of elementary row operations.
Equivalently, there exists an invertible $m\times m$ matrix $Q$ s.t. $QA$ is in its REF.

\begin{proof}
Suppose that $A$ is transformed into REF after \textcolor{red}{$k$
elementary row operations}. 

Let the corresponding elementary row matrices be $\textcolor{red}{E_1,E_2,\ldots,E_k}$.

Then $\textcolor{red}{(E_kE_{k-1}\cdots E_2E_1)}A$ is in REF.

We can take $Q=\textcolor{red}{E_kE_{k-1}\cdots E_2E_1}$ $\implies$ $QA$ is in row echelon form.
\end{proof}

\subsection{Example}

\subsubsection{Identity matrices examples}

The following are the three types of elementary row matrices obtained by
performing one elementary row operation on the identity matrix
\[
I_3=
\begin{pmatrix}
1&0&0\\
0&1&0\\
0&0&1
\end{pmatrix}.
\]


\[
E^{\mathrm{I}}_{1,2}=
\underbrace{
\begin{pmatrix}
0&1&0\\
1&0&0\\
0&0&1
\end{pmatrix}}_{\textcolor{red}{\text{$r_1 \leftrightarrow r_2$}}},
\qquad
E^{\mathrm{II}}_{\textcolor{red}{\lambda},3}=
\underbrace{
\begin{pmatrix}
1&0&0\\
0&1&0\\
0&0&\lambda
\end{pmatrix}}_{\textcolor{red}{\text{$r_3 \rightarrow \lambda r_3$}}},
\qquad
E^{\mathrm{III}}_{\textcolor{red}{\mu},2,3}=
\underbrace{
\begin{pmatrix}
1&0&0\\
0&1&0\\
0&\mu&1
\end{pmatrix}}_{\textcolor{red}{\text{$r_3 \rightarrow \mu r_2 + r_3$}}}.
\]

\subsubsection{Row operations $\leftrightarrow$ left multiplication}

Consider the 2 by 3 matrix
$A=
\begin{pmatrix}
1&2&3\\
3&4&5
\end{pmatrix}.
$
We can

\begin{enumerate}
    \item[(I)] Swap the first and second rows:
    \[
    \begin{pmatrix}
    3&4&5\\
    1&2&3
    \end{pmatrix}
    =
    \textcolor{red}{E^{\mathrm{I}}_{1,2}}A
    =
    \textcolor{red}{\begin{pmatrix}
    0&1\\
    1&0
    \end{pmatrix}}
    \begin{pmatrix}
    1&2&3\\
    3&4&5
    \end{pmatrix}.
    \]

    \item[(II)] Multiply the second row by $5$:
    \[
    \begin{pmatrix}
    1&2&3\\
    15&20&25
    \end{pmatrix}
    =
    \textcolor{red}{E^{\mathrm{II}}_{5,2}}A
    =
    \textcolor{red}{\begin{pmatrix}
    1&0\\
    0&5
    \end{pmatrix}}
    \begin{pmatrix}
    1&2&3\\
    3&4&5
    \end{pmatrix}.
    \]

    \item[(III)] Replace the second row by $-3r_1+r_2$:
    \[
    \begin{pmatrix}
    1&2&3\\
    0&-2&-4
    \end{pmatrix}
    =
    \textcolor{red}{E^{\mathrm{III}}_{-3,1,2}}A
    =
    \textcolor{red}{\begin{pmatrix}
    1&0\\
    -3&1
    \end{pmatrix}}
    \begin{pmatrix}
    1&2&3\\
    3&4&5
    \end{pmatrix}.
    \]
\end{enumerate}

\subsubsection{row echelon form (REF)}

1. The following matrices are in row echelon form (REF):
\[
\begin{pmatrix}
1&2&3&4&5\\
0&0&6&0&7\\
0&0&0&8&9
\end{pmatrix},
\quad
\begin{pmatrix}
1&2\\
0&3\\
0&0
\end{pmatrix},
\quad
\begin{pmatrix}
1&2&3\\
0&4&5\\
0&0&6
\end{pmatrix},
\quad
\begin{pmatrix}
1&2&3\\
0&0&4\\
0&0&0
\end{pmatrix},
\quad
\begin{pmatrix}
1&3&4&6\\
0&4&5&11\\
0&0&6&12
\end{pmatrix}.
\]

2. The following matrices are \emph{not} in row echelon form (REF):
\[
\begin{pmatrix}
1&2&3&4&5\\
0&0&0&6&7\\
0&0&8&0&9
\end{pmatrix},
\quad
\begin{pmatrix}
1&2\\
0&0\\
0&3
\end{pmatrix},
\quad
\begin{pmatrix}
1&2&3\\
0&0&4\\
0&0&1
\end{pmatrix},
\quad
\begin{pmatrix}
1&2&3\\
0&0&4\\
0&0&4
\end{pmatrix},
\quad
\begin{pmatrix}
1&3&4&6\\
0&4&5&11\\
0&6&0&12
\end{pmatrix}.
\]

\newpage

\subsubsection{compute REF}

Consider the matrix
\[
A=
\begin{pmatrix}
1&2&3&3&3\\
2&4&2&-4&2\\
0&0&1&2&1\\
3&6&3&-6&3
\end{pmatrix}.
\]

\begin{enumerate}
    \item[(a)] Transform $A$ into its row echelon form (REF).
    \item[(b)] Find the rank and nullity of $A$.
\end{enumerate}

\begin{solution}
(a) Perform the following elementary row operations:
\[
\begin{aligned}
R_2 &\leftarrow R_2-2R_1,\\
R_4 &\leftarrow R_4-3R_1,\\
R_3 &\leftarrow R_3+\frac14R_2,\\
R_4 &\leftarrow R_4-\frac32R_2.
\end{aligned}
\]
Hence,
\[
A
=
\begin{pmatrix}
1&2&3&3&3\\
0&0&4&-10&4\\
0&0&0&-\frac12&2\\
0&0&0&0&0
\end{pmatrix},
\]
which is in REF.

(b) Since the row echelon form has three nonzero rows, $\operatorname{rank}(A)=3$.

Because $A$ has $5$ columns, the Rank--Nullity Theorem gives $\operatorname{rank}(A)+\operatorname{nullity}(A)=\textcolor{red}{5}$.
Therefore, $\operatorname{nullity}(A)=5-3=\textcolor{red}{2}$.

\end{solution}

\newpage

\section{Column operations}

In analogy to elementary row operations, we can introduce
\emph{elementary column operations}. The corresponding theory is almost
identical, except that it has a more natural connection with the column
space and the rank of a matrix.

To begin with, we restate the definitions and theorems for row
operations in terms of columns.

There are three types of elementary column operations on a matrix $A$.

\begin{enumerate}
    \item[(I)] Interchange columns $i$ and $j$ of $A$.
    \item[(II)] Multiply the $i$-th column of $A$ by a nonzero scalar $\lambda$.
    \item[(III)] Add $\lambda$ times the $i$-th column to the $j$-th column.
\end{enumerate}

\subsection{Definition}

\textbf{Definition 6.2.1.}
An $(n\times n)$ elementary column matrix is a matrix obtained by
performing exactly one elementary column operation on the
$(n\times n)$ identity matrix $I_n$.

We denote the three types of elementary column matrices by
\[
E^{\mathrm{I}}_{i,j},\quad E^{\mathrm{II}}_{\lambda,i},\quad E^{\mathrm{III}}_{\lambda,i,j},
\]
respectively.

\medskip

\textbf{Definition 6.2.2.}
An $m\times n$ matrix $A$ is said to be in its
column echelon form (CEF) if all of the following conditions
are satisfied:

\begin{enumerate}
    \item Any column containing a nonzero entry is on the left of any
    zero columns.

    \item The first nonzero entry of each column occurs strictly below
    the first nonzero entry of the previous column.
\end{enumerate}

\subsection{Theorem}

\textbf{Theorem 6.2.1.}
Let $A$ be a matrix with $n$ columns. Then the matrices obtained by

\begin{enumerate}
    \item performing one of the column operations (I), (II), or (III) on $A$, and
    \item multiplying the corresponding elementary column matrix $E$ to the right of $A$, i.e., $AE$
\end{enumerate}

are the same.

\begin{proof}
This can be verified directly from the definition of matrix multiplication.
\end{proof}

\textbf{Theorem 6.2.2.}
Let $A$ be a matrix with $n$ columns. If $P$ is an invertible
$n\times n$ matrix, then we have
\[
\operatorname{rank}(AP)=\operatorname{rank}(A).
\]

In particular, multiplying an invertible matrix from the right preserves
the rank.
\begin{Remark}
Since elementary column operations correspond to right multiplication
by invertible elementary matrices, column operations preserve the rank.
\end{Remark}

\textbf{Theorem 6.2.3.}
Let $A$ be a matrix with $m$ columns
$\{\mathbf{c}_1,\ldots,\mathbf{c}_m\}$. Then

\begin{enumerate}
    \item $\operatorname{im}(A)=\operatorname{span}\{\mathbf{c}_1,\ldots,\mathbf{c}_m\}$.

    \item $\operatorname{rank}(A)$ equals the maximal number of linearly
    independent columns of $A$.

    \item $\operatorname{rank}(A)$ equals the number of nonzero columns
    when we transform $A$ into its column echelon form (CEF).
\end{enumerate}

\begin{proof}
We prove the three statements separately.

\medskip
\noindent
\textit{Proof of (1).}
By definition,
\begin{align*}
\operatorname{im}(A)
&=
\left\{
\textcolor{red}{A}\mathbf{x}:\mathbf{x}\in\mathbb{F}^m
\right\} \\
&=
\left\{
x_1\textcolor{red}{\mathbf{c}_1}+x_2\textcolor{red}{\mathbf{c}}_2+\cdots+x_m\textcolor{red}{\mathbf{c}_m}
:x_1,\ldots,x_m\in\mathbb{F}
\right\} \\
&=
\operatorname{span}
\{\textcolor{red}{\mathbf{c}_1,\mathbf{c}_2,\ldots,\mathbf{c}_m}\}.
\end{align*}

\medskip

\noindent

\textit{Proof of (2).}
Apply the sifting process to $\{\textcolor{red}{\mathbf{c}_1,\mathbf{c}_2,\ldots,\mathbf{c}_m}\}$.

We obtain a maximal linearly independent set formed by the columns

$\mathbf{c}_i$, which forms a basis of $\operatorname{im}(A)$ by (1).

Hence,the size of "sifited" set = $\dim\bigl(\operatorname{im}(A)\bigr)=\operatorname{rank}(A)$.

\medskip
\noindent
\textit{Proof of (3).}
A column \textcolor{red}{$\mathbf{c}_i$} is removed by the sifting process if and only if
the \textcolor{red}{$i$-th column} becomes a zero column after applying suitable elementary
column operations.

\end{proof}

\subsection{Example}

\subsubsection{column echelon form (CEF)}
Consider the matrix
\[
A=
\begin{pmatrix}
1&2&3\\
4&5&6
\end{pmatrix}
\in M_{2\times 3}(\mathbb{R}).
\]
Find $\operatorname{im}(A)$, a basis of $\operatorname{im}(A)$, and transform
$A$ into its column echelon form (CEF).

\begin{solution}
The image of $A$ is the span of its columns:
\[
\operatorname{im}(A)
=
\operatorname{span}
\left\{
\begin{pmatrix}1\\4\end{pmatrix},
\begin{pmatrix}2\\5\end{pmatrix},
\begin{pmatrix}3\\6\end{pmatrix}
\right\}.
\]

Applying the sifting process, we observe that
\[
\begin{pmatrix}3\\6\end{pmatrix}
=
\textcolor{red}{-1}\begin{pmatrix}1\\4\end{pmatrix}
+
\textcolor{red}{2}\begin{pmatrix}2\\5\end{pmatrix}.
\]
Hence, we obtain
\[
\left\{
\begin{pmatrix}1\\4\end{pmatrix},
\begin{pmatrix}2\\5\end{pmatrix}
\right\}
\]
forms a basis of $\operatorname{im}(A)$. In particular, $\operatorname{rank}(A)=2$.

Using elementary column operations,
\[
\begin{pmatrix}
1&2&3\\
4&5&6
\end{pmatrix}
\sim
\begin{pmatrix}
1&0&0\\
4&-3&0
\end{pmatrix},
\]
which is in column echelon form (CEF).
\end{solution}

\newpage

\section{Smith Normal Form}

\subsection{Definition}
\textbf{Definition 6.3.1.}
An $m\times n$ matrix $A$ is said to be in its
Smith normal form (SNF) if all of the following conditions
are satisfied:

\begin{enumerate}
    \item $A$ is both in row echelon form (REF) and column echelon form (CEF).
    \item All nonzero entries of $A$ are $1$.
\end{enumerate}

A typical matrix in Smith normal form is given below
\[
\begin{pmatrix}
I_r & 0_{r,n-r}\\
0_{m-r,r} & 0_{m-r,n-r}
\end{pmatrix},
\]
where $\mathbf{0}_{k,\ell}$ denotes the $k\times\ell$ zero matrix.

\subsection{Theorem}

\textbf{Theorem 6.3.1.} (QAP Theorem)
Let $A$ be an $m\times n$ matrix. Then there exist invertible matrices
$Q\in M_m(\mathbb{F})$ and $P\in M_n(\mathbb{F})$
s.t.
\[
QAP=
\begin{pmatrix}
I_r & 0_{r,n-r}\\
0_{m-r,r} & 0_{m-r,n-r}
\end{pmatrix}.
\]
Moreover, $r=\operatorname{rank}(A)$.

\begin{proof}
Given a matrix $A$, we first use \textcolor{red}{row operations} to transform $A$ into
its row echelon form (REF) $\Leftrightarrow$ \textcolor{red}{$\exists$ there  an invertible matrix
$Q$} s.t. $QA$ is in REF.

Next, we perform column operations on $QA$.

\begin{enumerate}
    \item First, use Type III column operations to "clean up" all
    entries to the right of the leading entries of each column.

    \item Second, use Type I column operations to move all newly
    obtained "zero columns" to the right.

    \item Third, use Type II column operations to "rescale" all nonzero
    entries to $1$.
\end{enumerate}

Hence, the resulting matrix is in Smith normal form. Since column
operations correspond to right multiplication by elementary matrices,
there $\exists$ an invertible matrix $P$ s.t.
\[
\textcolor{red}{QA}P
=
\underbrace{
\begin{pmatrix}
I_r & 0_{r,n-r}\\
0_{m-r,r} & 0_{m-r,n-r}
\end{pmatrix}.}_{\textcolor{red}{\text{SNF}}}
\]

Moreover, multiplication by invertible matrices preserves rank.
Therefore,
\[
\operatorname{rank}(A)
=
\operatorname{rank}(QAP)
=
\operatorname{rank}
\begin{pmatrix}
I_r & 0_{r,n-r}\\
0_{m-r,r} & 0_{m-r,n-r}
\end{pmatrix}
=
r.
\]

Thus, $r=\operatorname{rank}(A)$.
\end{proof}

\textbf{Theorem 6.3.2.}
Let $A$ be an $m\times n$ matrix. Then $\operatorname{rank}(A)=\operatorname{rank}(A^t)$.

\begin{Remark}
In some textbooks, this theorem is referred to as
'row rank = column rank'.
\end{Remark}

\begin{proof}
Given $A$, by the \textcolor{red}{QAP Theorem}, there exist invertible matrices
$P$ and $Q$ s.t.
\[
QAP=
\begin{pmatrix}
\textcolor{red}{I_r} & 0\\
0 & 0
\end{pmatrix}.
\]

Taking transpose, we have
\[
P^tA^tQ^t
=
(QAP)^t
=
\begin{pmatrix}
\textcolor{red}{I_r} & 0\\
0 & 0
\end{pmatrix}.
\]

Therefore,
\[
\operatorname{rank}(A)
=
\operatorname{rank}(QAP)
=
\textcolor{red}{r}
=
\operatorname{rank}(P^tA^tQ^t)
\leq
\operatorname{rank}(A^t),
\]
where the first equality follows since $P$ and $Q$ are invertible.

Since $P^t$ and $Q^t$ are also invertible, the same argument applied
in the reverse direction gives
\[
\operatorname{rank}(A^t)
\leq
\operatorname{rank}(A).
\]

Hence, $\operatorname{rank}(A)=\operatorname{rank}(A^t)$.
\end{proof}

\textbf{Theorem 6.3.3.}
Let $A$ be an $m\times n$ matrix. Then the following quantities are equal:

\begin{enumerate}
    \item the rank of $A$;
    \item the maximal number of linearly independent columns of $A$;
    \item the number of nonzero columns in CEF;
    \item the rank of $A^t$;
    \item the maximal number of linearly independent rows of $A$;
    \item the number of nonzero rows in REF.
\end{enumerate}

\begin{proof}
By the previous results,
\begin{align*}
\operatorname{rank}(A)
&= \operatorname{rank}(A^t) \\
&= \text{the maximal number of linearly independent columns of } A^t \\
&= \text{the number of nonzero columns in CEF} \\
&= \text{the maximal number of linearly independent} \textcolor{red}{\text{ rows of A}} \\
&= \text{the number of nonzero} \textcolor{red}{\text{ rows in REF}}.
\end{align*}
\end{proof}

\end{document}